% GENERATED FILE. Edit canonical lessons, then run npm run generate:books.
% canonical-source-hash: fnv1a64:2810b703278aad1b
% canonical-lessons: ZH-C510-waipo, ZH-C510-sunzi, ZH-C510-sunnu, ZH-C510-tengai, ZH-C510-hong

\chapter{Grandparents}
\label{ch:zh-grandparents}
\hlchaptermodality{510}

\begin{chapteropening}
\textbf{By the end of this chapter:} \emph{I can talk about my grandma on my mother's side, grandsons and granddaughters, doting and soothing a child.}

Complete the last lesson of chapter 510: I can talk about my grandma on my mother's side, grandsons and granddaughters, doting and soothing a child.
\end{chapteropening}

\section[wàipó]{wàipó — grandma on the mother's side}

\label{lesson:ZH-C510-waipo}



\begin{quote}
Before the new one: say the Mandarin for one's mother tongue, then the Mandarin for a stepmother.
\end{quote}

\subsection*{You'll want to know: wàipó}
\emph{wàipó} — \textquotedblleft{}grandma on the mother's side\textquotedblright{}.

\begin{quote}
\emph{wǒ zài wàipó jiā zhù le yí ge yuè} — I stayed at my grandma's for a month.
\end{quote}

Its characters come later; for now it is a word for the ear.

\begin{cousinweb}
\emph{wàipó} is built from \zh{外} 'outside' and \emph{pó} 'an old woman', both inherited from Middle Chinese, \emph{ngwajH} and \emph{ba} (Guangyun, in Baxter's transcription), which glosses the second 'a name for an old mother'.
\end{cousinweb}

\begin{grammarlens}[title={What you've built}]
One more word for saying how it was.
\end{grammarlens}

\subsection*{Your turn}
\begin{itemize}
\raggedright
  \item \emph{Say it:} \emph{wàipó}
  \item \emph{Say it:} \emph{wàipó}, once more
  \item \emph{Say it:} \emph{wǒ zài wàipó jiā zhù le yí ge yuè} — I stayed at my grandma's for a month.
  \item \emph{Recall:} say the Mandarin for one's mother tongue, then the Mandarin for a stepmother, then say \emph{wàipó} again
\end{itemize}

\begin{tcolorbox}[breakable,colback=teal!4,colframe=teal!35!black,title={Before you move on}]
What does wàipó mean? (\textquotedblleft{}Grandma on the mother's side\textquotedblright{}.) Say it once more.
\end{tcolorbox}
\section[sūnzi]{sūnzi — grandson}

\label{lesson:ZH-C510-sunzi}



\begin{quote}
Before the new one: say the Mandarin for a stepmother, then the Mandarin for grandma on the mother's side.
\end{quote}

\subsection*{You'll want to know: sūnzi}
\emph{sūnzi} — \textquotedblleft{}grandson\textquotedblright{}.

\begin{quote}
\emph{yéye zuì xǐhuan tā de sūnzi} — Grandpa loves his grandson most of all.
\end{quote}

Its characters come later; for now it is a word for the ear.

\begin{cousinweb}
\emph{sūnzi} is built from \emph{sūn} 'a grandchild' and \zh{子} 'a son', both inherited from Middle Chinese, \emph{swon} and \emph{tsiX} (Guangyun, in Baxter's transcription), which glosses the first, after an older dictionary, 'the son of a son'.
\end{cousinweb}

\begin{grammarlens}[title={What you've built}]
One more word for saying how it was.
\end{grammarlens}

\subsection*{Your turn}
\begin{itemize}
\raggedright
  \item \emph{Say it:} \emph{sūnzi}
  \item \emph{Say it:} \emph{sūnzi}, once more
  \item \emph{Say it:} \emph{yéye zuì xǐhuan tā de sūnzi} — Grandpa loves his grandson most of all.
  \item \emph{Recall:} say the Mandarin for a stepmother, then the Mandarin for grandma on the mother's side, then say \emph{sūnzi} again
\end{itemize}

\begin{tcolorbox}[breakable,colback=teal!4,colframe=teal!35!black,title={Before you move on}]
What does sūnzi mean? (\textquotedblleft{}Grandson\textquotedblright{}.) Say it once more.
\end{tcolorbox}
\section[sūnnǚ]{sūnnǚ — granddaughter}

\label{lesson:ZH-C510-sunnu}



\begin{quote}
Before the new one: say the Mandarin for grandma on the mother's side, then the Mandarin for a grandson.
\end{quote}

\subsection*{You'll want to know: sūnnǚ}
\emph{sūnnǚ} — \textquotedblleft{}granddaughter\textquotedblright{}.

\begin{quote}
\emph{nǎinai měi tiān sòng sūnnǚ qù shàngxué} — Grandma takes her granddaughter to school every day.
\end{quote}

Its characters come later; for now it is a word for the ear.

\begin{cousinweb}
\emph{sūnnǚ} is built from \emph{sūn} 'a grandchild' and \zh{女} 'a daughter', both inherited from Middle Chinese, \emph{swon} and \emph{nrjoX} (Guangyun, in Baxter's transcription).
\end{cousinweb}

\begin{grammarlens}[title={What you've built}]
One more word for saying how it was.
\end{grammarlens}

\subsection*{Your turn}
\begin{itemize}
\raggedright
  \item \emph{Say it:} \emph{sūnnǚ}
  \item \emph{Say it:} \emph{sūnnǚ}, once more
  \item \emph{Say it:} \emph{nǎinai měi tiān sòng sūnnǚ qù shàngxué} — Grandma takes her granddaughter to school every day.
  \item \emph{Recall:} say the Mandarin for grandma on the mother's side, then the Mandarin for a grandson, then say \emph{sūnnǚ} again
\end{itemize}

\begin{tcolorbox}[breakable,colback=teal!4,colframe=teal!35!black,title={Before you move on}]
What does sūnnǚ mean? (\textquotedblleft{}Granddaughter\textquotedblright{}.) Say it once more.
\end{tcolorbox}
\section[\texorpdfstring{\zh{疼爱} (téng'ài)}{疼爱 (téng'ài)}]{\zh{疼爱} (téng'ài) — to dote on}

\label{lesson:ZH-C510-tengai}



\begin{quote}
Before the new one: say the Mandarin for a grandson, then the Mandarin for a granddaughter.
\end{quote}

\subsection*{You'll want to know: \zh{疼爱}}
\textbf{\zh{疼爱}} — \emph{téng'ài} — \textquotedblleft{}to dote on\textquotedblright{}.

\begin{quote}
\emph{wàipó hěn téng'ài wǒ} — Grandma really dotes on me.
\end{quote}

\begin{cousinweb}
\textbf{\zh{疼爱}} (\emph{téng'ài}) is built from \zh{疼} 'to ache' and \emph{ài} 'to love', both inherited from Middle Chinese, \emph{downg} and \emph{'ojH} (Guangyun, in Baxter's transcription), which glosses the first 'pain' and the second 'to pity'; the sense 'to love dearly' is not given for the first there.
\end{cousinweb}

\begin{grammarlens}[title={What you've built}]
One more word for saying how it was.
\end{grammarlens}

\subsection*{Your turn}
\begin{itemize}
\raggedright
  \item \emph{Say it:} \emph{téng'ài}
  \item \emph{Say it:} \emph{téng'ài}, once more
  \item \emph{Say it:} \emph{wàipó hěn téng'ài wǒ} — Grandma really dotes on me.
  \item \emph{Recall:} say the Mandarin for a grandson, then the Mandarin for a granddaughter, then say \emph{téng'ài} again
\end{itemize}

\begin{tcolorbox}[breakable,colback=teal!4,colframe=teal!35!black,title={Before you move on}]
What does \zh{疼爱} mean? (\textquotedblleft{}To dote on\textquotedblright{}.) Say it once more.
\end{tcolorbox}
\section[hǒng]{hǒng — to soothe a child; to coax}

\label{lesson:ZH-C510-hong}



\begin{quote}
Before the new one: say the Mandarin for a granddaughter, then the Mandarin for to dote on.
\end{quote}

\subsection*{You'll want to know: hǒng}
\emph{hǒng} — \textquotedblleft{}to soothe a child; to coax\textquotedblright{}.

\begin{quote}
\emph{wàipó zuì huì hǒng háizi} — Grandma is the best at soothing children.
\end{quote}

Its characters come later; for now it is a word for the ear.

\begin{cousinweb}
\emph{hǒng} is written with a character the Guangyun lists only as Middle Chinese \emph{huwngH} (Guangyun, in Baxter's transcription), glossed 'the sound of singing', which would regularly give hòng. The sense 'to soothe a child' is not given there and the third tone does not fit, so how the word came is not settled.
\end{cousinweb}

\begin{grammarlens}[title={What you've built}]
One more word for saying how it was.
\end{grammarlens}

\subsection*{Your turn}
\begin{itemize}
\raggedright
  \item \emph{Say it:} \emph{hǒng}
  \item \emph{Say it:} \emph{hǒng}, once more
  \item \emph{Say it:} \emph{wàipó zuì huì hǒng háizi} — Grandma is the best at soothing children.
  \item \emph{Recall:} say the Mandarin for a granddaughter, then the Mandarin for to dote on, then say \emph{hǒng} again
\end{itemize}

\begin{tcolorbox}[breakable,colback=teal!4,colframe=teal!35!black,title={Before you move on}]
What does hǒng mean? (\textquotedblleft{}To soothe a child; to coax\textquotedblright{}.) Say it once more.
\end{tcolorbox}
